% ----- ----- ----- ----- ----- ----- ----- ----- ----- ----- 
% Preamble
% ----- ----- ----- ----- ----- ----- ----- ----- ----- ----- 
\documentclass[a4paper,9pt]{article}
\usepackage[utf8]{inputenc}
\usepackage[english]{babel}
\usepackage{amsmath}
\usepackage{authblk}
\usepackage{background}
\usepackage{enumitem}
\usepackage{fancyhdr}
\usepackage{geometry}
\usepackage{graphicx}
\usepackage[hidelinks]{hyperref}
\usepackage{indentfirst}
\usepackage{multicol}
\usepackage{tabularx}
\usepackage[most]{tcolorbox}
\usepackage{tgheros}
\usepackage[explicit]{titlesec}
\usepackage{xcolor}
% ----- ----- ----- ----- ----- 
% Column Spacing
% ----- ----- ----- ----- ----- 
\setlength{\columnseprule}{0pt}
\setlength{\columnsep}{10.0pt}
\newcolumntype{L}[1]{>{\raggedright\let\newline\\\arraybackslash\hspace{0pt}}m{#1}}
\newcolumntype{C}[1]{>{\centering\let\newline\\\arraybackslash\hspace{0pt}}m{#1}}
\newcolumntype{R}[1]{>{\raggedleft\let\newline\\\arraybackslash\hspace{0pt}}m{#1}}
% ----- ----- ----- ----- ----- 
% Custom Colors
% ----- ----- ----- ----- ----- 
\definecolor{myDColor}{HTML}{000000} % Darker Color
\definecolor{myLColor}{HTML}{CFB87C} % Lighter Color
\definecolor{LinkColor}{HTML}{CFB87C} % Link Color
% ----- ----- ----- ----- ----- 
% Custom Title
% ----- ----- ----- ----- ----- 
\makeatletter
\renewcommand*{\maketitle}{%
\noindent
\begin{minipage}{\textwidth}
\begin{tikzpicture}
\node[rectangle,rounded corners=10pt,inner sep=7.5pt,fill=myDColor,text width= 0.975\textwidth, align=center] 
{\color{white}\Huge \@title};
\end{tikzpicture}
\end{minipage}
\hfill
\bigskip\bigskip
}%
\makeatother
% ----- ----- ----- ----- ----- 
% Custom Section
% ----- ----- ----- ----- ----- 
\newcommand*\sectionlabel{}
\titleformat{\section}
{\gdef\sectionlabel{}
\normalfont\sffamily\Large\bfseries\scshape}
{\gdef\sectionlabel{\thesection\ }}{0pt}{
\noindent
\begin{tikzpicture}
\node[rectangle,rounded corners=6pt,inner sep=4pt,fill=myLColor,text width= 0.975\columnwidth, align=center] 
{\color{white}\sectionlabel#1};
\end{tikzpicture}}
\titlespacing*{\section}{0pt}{5pt}{5pt}
% ----- ----- ----- ----- ----- 
% Custom Header & Footer / Geometry
% ----- ----- ----- ----- ----- 
\geometry{a4paper,left=10mm,right=10mm,top=15mm,bottom=15mm}
\linespread{1} % Distance between lines
\pagestyle{fancy}
\fancyhead[L]{Taylor Larrechea}
\fancyhead[C]{Linked List Assignment Interview}
\fancyhead[R]{University of Colorado}
\fancyfoot[L]{CSPB 2270 Data Structures}
\fancyfoot[R]{Linked List AI}
\backgroundsetup{
    scale = 1,
    angle = 0,
    opacity = 0.1,
    contents = {
    \includegraphics[scale = 0.25, keepaspectratio]{Figures/CU Seal.png}
    }
}
\newtcolorbox{highlight}[1][]{%
    enhanced,
    skin first=enhanced,
    skin middle=enhanced,
    skin last=enhanced,
    before upper={\parindent15pt},
    breakable,
    boxrule = 0pt,
    frame hidden,
    borderline west = {4pt}{0pt}{myDColor},
    colback = myLColor!10,
    coltitle = myLColor!5,
    sharp corners,
    rounded corners = southeast,
    arc is angular,
    arc = 3mm,
    attach boxed title to top left,
    boxed title style = {%
        enhanced,
        colback = myDColor,
        colframe = myDColor,
        top = 0pt,
        bottom = 0pt,
        sharp corners,
        rounded corners = northeast,
        arc is angular,
        arc = 2mm,
        rightrule = 0pt,
        bottomrule = 0pt,
        toprule = 0pt,
    },
    title = {\bfseries\large #1}, 
    overlay unbroken={%
        \path[fill = tcbcolback!80!black] 
            ([yshift = 3mm]interior.south east) -- ++(-0.4,-0.1) -- ++(0.1,-0.2);
    },
    overlay first = {%
        \path[fill = tcbcolback!80!black] 
            ([yshift = 3mm]interior.south east) -- ++(-0.4,-0.1) -- ++(0.1,-0.2);
    },
    overlay middle={%
        \path[fill = tcbcolback!80!black] 
            ([yshift = -3mm]interior.north east) -- ++(-0.4,0.1) -- ++(0.1,0.2);
        \path[fill = tcbcolback!80!black] 
            ([yshift = 3mm]interior.south east) -- ++(-0.4,-0.1) -- ++(0.1,-0.2);
    },
    overlay last={%
        \path[fill = tcbcolback!80!black] 
            ([yshift = -3mm]interior.north east) -- ++(-0.4,0.1) -- ++(0.1,0.2);
        \path[fill = tcbcolback!80!black] 
            ([yshift = 3mm]interior.south east) -- ++(-0.4,-0.1) -- ++(0.1,-0.2);
    },
    extras middle and last = { rounded corners = northeast }
}
\newcommand{\horizontalline}{\noindent \rule{\textwidth}{0.5pt}\par}
% ----- ----- ----- ----- ----- Title
\title{Linked List Assignment Interview Notes}
% ----- ----- ----- ----- ----- Begin Document
\begin{document}

\maketitle

\section*{Vector Overview}

\subsection*{Generic Overview}

A vector is a dynamic array that allows storing and accessing elements of the same data type. It is a 
container class
provided by the Standard Template Library (STL) in C++. Vectors provide several advantages over traditional 
arrays, such
as automatic memory management, resizable capacity, and built-in functions for efficient manipulation.

Internally, a vector is implemented as a contiguous block of memory that stores the elements. It dynamically 
manages its
size, expanding or shrinking as needed. Elements are accessed using zero-based indexing, similar to arrays. 
The vector
class provides various member functions and operators for performing common operations like adding elements 
at the end,
inserting elements at specific positions, removing elements, and accessing elements by index. Vectors offer 
efficient
random access, constant time complexity for accessing elements, and amortized constant time complexity for 
adding or 
removing elements at the end.

\subsection*{How Are Vectors Different From Arrays?}

\noindent Vectors and arrays are both used to store collections of elements in C++, but vary in some areas:

\begin{itemize}
    \item Arrays are fixed-size collections with a static size determined at compile-time, while vectors are 
dynamic arrays
    with a flexible and resizable size.
    \item Arrays require manual memory management and cannot be easily resized, whereas vectors handle memory 
allocation and
    deallocation automatically, allowing for dynamic resizing.
    \item Arrays provide direct and efficient access to elements using zero-based indexing, while vectors 
offer convenient
    member functions like `push\_back()' and `pop\_back()' for adding and removing elements.
    \item Arrays have a fixed memory layout, resulting in slightly faster element access, while vectors 
provide more convenience
    and flexibility.
    \item Arrays cannot be easily copied or assigned to another array, while vectors can be easily copied of 
assigned to other
    vectors.
\end{itemize}

\noindent In summary, arrays are suitable when the size is known and fixed, and direct memory control is 
needed. Vectors, on the other hand,
are preferred when the size needs to be dynamic, convenient operations are required, and automatic memory 
management is desired.

\subsection*{What Is The Algorithm For Adding And Removing Elements From A Vector?}

\noindent To add elements to a vector in C++, we use:

\begin{itemize}
    \item To add an element at the end of the vector, we use the `push\_back()' function, which appends the 
element to the end of
    the vector.
    \item To insert an element at a specific position within the vector, we use `insert()' function. This 
function takes an iterator
    pointing to the position and the element to be inserted.
\end{itemize}

\noindent To remove elements from a vector in C++, we use:

\begin{itemize}
    \item To remove the last element of the vector, we use the `pop\_back()' function, which removes the 
element from the end of the vector.
    \item If you want to remove an element from a specific position within the vector, we use the `erase()' 
function. It takes an iterator
    pointing to the position and removes the element at that position.
\end{itemize}

\noindent These operations provide a convenient and efficient way to add and remove elements from a vector, 
maintaining the integrity and
resizing the vector dynamically as needed.

\subsection*{What Are The Benefits Of Using A Vector Over An Array?}

\noindent The benefits of using a vector over an array are:

\begin{itemize}
    \item Dynamic Size: Vectors provide dynamic resizing, allowing you to add or remove elements at runtime. 
Unlike arrays, which have a fixed
    size, vectors automatically handle memory allocation.
    \item Convenience: Vectors come with built-in functions like `push\_back()', `pop\_back()', and 
`insert()', making it easier to add, remove,
    or insert elements without manual management. Vectors also provide member functions like `size()' and 
`empty()' to retrieve information about
    the number of elements and check if the vector is empty.
    \item Bounds Checking: Vectors perform bounds checking on element access, ensuring that you do not go out 
of bounds. This helps prevent accessing
    memory beyond the allocated size and reduces the risk of errors and crashes that can occur with arrays.
    \item Iteration Support: Vectors can be easily iterated using range-based for loops or iterators, 
providing a convenient way to traverse and
    manipulate elements.
    \item Copying and Assignment: Vectors can be easily copied and assigned to other vectors using the copy 
constructor and assignment operator, allowing
    for efficient and convenient data handling.
\end{itemize}

\noindent Overall, vectors offer the flexibility of dynamic resizing, convenient functions, bounds checking, 
iteration support, and ease of copying,
making them a preferred choice over arrays in many scenarios.

\subsection*{What Is The Basic Algorithm For Iterating Through A Vector?}

The simple algorithm for iterating through the elements of a vector is to use a for-loop. Starting the index 
at 0, and going up to the size of
the vector, we can use built in functions like `at()' to indicate which index we are wanting to examine / 
manipulate. Using this simple algorithm, we check
for certain conditions within our vector, look for specific elements, and many other things as well.

\section*{Linked List Overview}

\subsection*{Generic Overview}

A linked list is a data structure commonly used in object-oriented programming that consists of nodes 
connected through 
pointers. Each node contains a value and a pointer to the next node in the list, except for the last node 
which points 
to `nullptr'. The linked list does not have a specific `head' or `tail' node, but rather, the first node in 
the list is 
often referred to as the `head' and the last node is commonly known as the `tail'. The nodes in between the 
`head' and 
`tail' are linked in a sequential manner. The `head' node serves as the starting point for traversing the 
list, and each 
node points to the next node, allowing efficient insertion and removal operations.

\subsection*{Difference Between An Array And Vector}

Linked lists differ from arrays and vectors in several ways. One significant distinction is that the nodes in 
a linked list, 
which correspond to elements in an array or vector, contain pointers that indicate the next node in the list. 
In contrast, 
elements in arrays and vectors do not have pointers pointing to the next element. Unlike arrays and vectors, 
linked lists 
cannot be accessed using an index value. However, similar to vectors, linked lists have the potential to 
dynamically resize 
after compilation. It's important to note that the mechanism of size change in a linked list is different 
from that of vectors.
Some other key differences between linked lists and vectors and arrays are:

\begin{itemize}
    \item Memory Allocation: Linked lists dynamically allocate memory for each node as it is needed. This 
allows for efficient memory
    usage, as nodes can be allocated and deallocated independently. In contrast, arrays and vectors allocate 
a fixed block of memory
    upfront, regardless of the number of elements, which can lead to potential wastage of insufficient 
memory.
    \item Insertion and Deletion: Linked lists excel in insertion and deletion operations, especially in the 
middle of the list. Inserting
    or removing a node in a linked list only requires updating the pointers, whereas in arrays and vectors, 
these operations may involve
    shifting elements, resulting in less efficiency.
    \item Random Access: Arrays and vectors support direct and efficient random access using index-based 
access. You can access any element
    in constant time with their respective indices. Linked lists, on the other hand, do not provide direct 
index-based access, requiring traversal
    from the head to the desired node, resulting in linear time complexity.
    \item Memory Overhead: Linked lists have a higher memory overhead compared to arrays and vectors due to 
the additional memory required for storing
    the pointers. Each node in a linked list needs to store a pointer to the next node, increasing the 
overall memory usage.
    \item Iteration efficiency: Arrays and vectors offer efficient iteration using loops with index-based 
access. Since the elements are stored in
    contiguous memory locations, sequential access is faster. Linked lists , on the other hand, require 
traversing each node sequentially, resulting in
    slower iteration.
    \item Memory Fragmention: Linked lists are less prone to memory fragmentation compared to arrays and 
vectors. As nodes are dynamically allocated, they
    cam be scattered in memory, reducing the likelihood of large blocks of unallocated space.
\end{itemize}

\subsection*{What Is A Node In A Linked List?}

A node in a linked list is analogous to an element in an array or vector. In a linked list, a node contains 
data and a reference (or pointer) to the next node in the list.

\begin{itemize}
    \item Data Component: Represents the value or information stored in the node
    \item Reference Component: Points to the next node in the linked list with the use of a pointer
\end{itemize}

\noindent A node is analogous to a link in a chain. Each addition of a node will elongate the total length of 
the chain. 

\subsection*{Where Are Elements Accessed And Modified?}

The main examples of accessing and modifying elements in a linked list are the following:

\begin{itemize}
    \item Accessing the Value of a Node: We use the `data' member to access the value stored in a node of a 
linked list.
    \item Modifying the Value of a Node: To modify the value of a node, we use the syntax: `nodePtr-$>$data = 
newValue', or some variation like this.
    \item Traversing the Linked List: To traverse through the linked list, we use a while loop to do so.
    \item Inserting a New Node: We use the `Insert' functions defined in the program to insert nodes into the 
list.
    \item Removing a Node From the List: We use the 'Remove' functions defined in the program to remove nodes 
from the list.
\end{itemize}

\begin{highlight}[Accessing \& Modifying Elements]
    We access and modify elements in a linked list with function calls and the like below:

    \horizontalline

    \begin{verbatim}
    \\ Accessing the values of a node and traversing a linked list 
    \\ like in the Contains function below:
    bool LinkedList::Contains(int data){
        shared_ptr<node> currPtr = top_ptr_;
        bool ret = false;
        while (currPtr != nullptr) {
        if (currPtr->data == data) {
            ret = true;
        }
        else {}
        currPtr = currPtr->next;
        }
        return ret;
    }

    \\ Modifying the value of a node like in the AppendData function below:
    void LinkedList::AppendData(int data){
        shared_ptr<node> nxtPtr(new node);
        nxtPtr->data = data;
        nxtPtr->next = nullptr;
        if (top_ptr_ == nullptr) {
        SetTop(nxtPtr);
        }
        else {
        shared_ptr<node> currPtr = top_ptr_;
        while (currPtr->next != nullptr) {
            currPtr = currPtr->next;
        }
        currPtr->next = nxtPtr;
        }
    }

    \\ Inserting a new node into the linked list like in the Insert function below:
    void LinkedList::Insert(int offset, shared_ptr<node> new_node){
        shared_ptr<node> currPtr = top_ptr_;
        vector<shared_ptr<node>> ptrs;
        while (currPtr != nullptr) {
        ptrs.push_back(currPtr);
        currPtr = currPtr->next;
        }
        new_node->data;
        new_node->next = nullptr;
        if (offset == 0) {
        new_node->next = top_ptr_;
        SetTop(new_node);
        }
        else {
        currPtr = top_ptr_;
        for (int i = 0; i < ptrs.size(); i++) {
            if (i == offset - 1) {
            new_node->next = currPtr->next;
            currPtr->next = new_node;
            }
            else {}
            currPtr = currPtr->next;
            }
        }
    }

    \\ Removing a node from the linked list like in the Remove function below:
    void LinkedList::Remove(int offset){
        shared_ptr<node> currPtr = top_ptr_;
        vector<shared_ptr<node>> ptrs;
        while (currPtr != nullptr) {
        ptrs.push_back(currPtr);
        currPtr = currPtr->next;
        }
        if (offset == 0) {
        SetTop(top_ptr_->next);
        }
        else {
        currPtr = top_ptr_;
        for (int i = 0; i < ptrs.size(); i++) {
            if (i == offset - 1) {
            currPtr->next = currPtr->next->next;
            break;
            }
            else {}
            currPtr = currPtr->next;
            }
        }
    }
    \end{verbatim}

    \horizontalline
\end{highlight}

\subsection*{Initializing A Linked List}

When the linked list is constructed, we initialize the data member `next' to be null. To initialize the list 
with values, we call
the `InitNode' function that is defined in the CPP file from the project. This algorithm can be seen below:

\begin{highlight}[Value Initialization]

    \begin{verbatim}
    /*  InitNode - This function takes in an integer data type and assigns that value to a "node" 
    data type's "data" member
    *   Input:
    *     data - This is an integer data type that is later assigned to a "node" data type's "data" 
    member
    *   Algorithm:
    *     shared_ptr<node> ret(new node); - This line creates a smart pointer of data type "node" 
    named "ret"
    *   Output:
    *     ret - After setting "ret"'s data member to the input parameter "data"
    *     The "next" data member of the "node" object is initially set to nullptr, indicating the 
    absence of a next node
    */
    shared_ptr<node> LinkedList::InitNode(int data){
        shared_ptr<node> ret(new node);
        ret->data = data;
        return ret;
    }
    \end{verbatim}

\end{highlight}

\subsection*{Adding Elements To A Linked List}

Below is the algorithm for adding elements to a linked list found in this class.

\begin{highlight}[Insert Algorithm]
    \begin{verbatim}
    /*  Insert - This function inserts a new node into a linked list at the specified offset
    *   Input:
    *     offset - An integer that is the offset at which the new node should be inserted
    *     new_node - A smart pointer of object node that is to be inserted into the linked list
    *   Algorithm:
    *     First we create a new smart pointer called "currPtr" and set it equal to the top of 
    the linked list
    *     After that we create a vector called "ptrs" to store the smart pointers to each node 
    in the linked list
    *     Then we traverse the linked list and add each node's smart pointer to the "ptrs" vector
    *     After that, we check if the offset is zero:
    *       If it is zero, set the next pointer of the "new_node" node to the current top_ptr_
    *       Set the top_ptr_ to the "new_node" node, making it the new first node in the 
    linked list
    *     If the offset is not zero:
    *       Reset "currPtr" to the top_ptr_ and iterate through the "ptrs" vector
    *       When the current index matches the offset - 1, set the next pointer of "new_node" to 
    currPtr->next
    *       Update the next pointer of "currPtr" to "insertPtr", inserting it at the desired 
    offset
    *   Output:
    *     This function does not return a value
    */
    void LinkedList::Insert(int offset, shared_ptr<node> new_node){
        shared_ptr<node> currPtr = top_ptr_;
        vector<shared_ptr<node>> ptrs;
        while (currPtr != nullptr) {
        ptrs.push_back(currPtr);
        currPtr = currPtr->next;
        }
        new_node->data;
        new_node->next = nullptr;
        if (offset == 0) {
        new_node->next = top_ptr_;
        SetTop(new_node);
        }
        else {
        currPtr = top_ptr_;
        for (int i = 0; i < ptrs.size(); i++) {
            if (i == offset - 1) {
            new_node->next = currPtr->next;
            currPtr->next = new_node;
            }
            else {}
            currPtr = currPtr->next;
            }
        }
    }
    \end{verbatim}
\end{highlight}

\subsection*{Removing Elements From A Linked List}

Below is the algorithm for removing elements from a linked list found in this class.

\begin{highlight}[Remove Algorithm]
    \begin{verbatim}
    /*  Remove - This function removes a node from a linked list at the specified offset
    *   Input:
    *     offset - An integer that is the offset of the node to be removed from our linked list
    *   Algorithm:
    *     First we create a new smart pointer called "currPtr" and set it equal to the top of the 
    linked list
    *     Then we create a vector called "ptrs" to store the smart pointers to each node in the 
    linked list
    *     After that we iterate through the linked list and add each node's smart pointer to the 
    "ptrs" vector
    *     Then we heck if the offset is zero.
    *       If it is zero, set the next pointer of the top node to the node after it
    *     Otherwise, reset "currPtr" to the top node and iterate through the "ptrs" vector
    *       When the current index matches the offset - 1, set the next pointer of "currPtr" to 
    the node after it
    *       Break out of the for loop to avoid segmentation faults
    *     Output:
    *       This function does not return a value.
    */
    void LinkedList::Remove(int offset){
        shared_ptr<node> currPtr = top_ptr_;
        vector<shared_ptr<node>> ptrs;
        while (currPtr != nullptr) {
        ptrs.push_back(currPtr);
        currPtr = currPtr->next;
        }
        if (offset == 0) {
        SetTop(top_ptr_->next);
        }
        else {
        currPtr = top_ptr_;
        for (int i = 0; i < ptrs.size(); i++) {
            if (i == offset - 1) {
            currPtr->next = currPtr->next->next;
            break;
            }
            else {}
            currPtr = currPtr->next;
            }
        }
    }
    \end{verbatim}
\end{highlight}

\subsection*{Finding Elements In A Linked List}

Below is the algorithm for finding elements in a linked list found in this class:

\begin{highlight}[Search Algorithm]
    \begin{verbatim}
    /*  Contains - This function determines if an element is present inside a linked list
    *   Input:
    *     data - This is an integer that is being searched for in our linked list
    *   Algorithm:
    *     First, create a smart pointer called "currPtr" and initialize it to the head of 
    our linked list
    *     Initialize the return value "ret" to false, so that if a value is not found we do 
    not need to change it to false
    *     Traverse through the linked list with the use of a while loop and check if the 
    current nodes "data" member matches
    *     the input parameter "data"
    *       If the two are a match, change the value of "ret" to true, otherwise, leave it 
    as false
    *   Output:
    *     ret - Boolean value indicating if a value was found in the linked list
    */
    bool LinkedList::Contains(int data){
        shared_ptr<node> currPtr = top_ptr_;
        bool ret = false;
        while (currPtr != nullptr) {
        if (currPtr->data == data) {
            ret = true;
        }
        else {}
        currPtr = currPtr->next;
        }
        return ret;
    }
    \end{verbatim}
\end{highlight}

\subsection*{Returning The Size Of A Linked List}

Below is the algorithm for determining the size of a linked list found in this class:

\begin{highlight}[Size Algorithm]
    \begin{verbatim}
    /*  Size - This function counts how many nodes are present in our linked list
    *   Algorithm:
    *     We first create a smart pointer called "currPtr" and set it equal to our "top_ptr_" 
    in our linked list
    *     We initialize the return value "ret" to zero
    *     Then we enter a while loop that increments "ret" by one until it runs into the tail 
    of the list
    *     Advance the current pointer object using "currPtr = currPtr->next;"
    *   Output:
    *     ret - Integer value that represents the number of nodes that are present in our 
    linked list
    */
    int LinkedList::Size(){
        shared_ptr<node> currPtr = top_ptr_;
        int ret = 0;
        while (currPtr != nullptr) {
        ret++;
        currPtr = currPtr->next;
        }
        return ret;
    }
    \end{verbatim}
\end{highlight}

\subsection*{Reversing A Linked List}

The easiest way to reverse a linked list is to have a doubly linked list. With a singly linked list, the 
simplest way to
reverse the linked list is to create three separate nodes: `previous', `current', and `next'. Once these 
nodes have been
created, we traverse through the linked list and update their values accordingly. This algorithm can be seen 
below:

\begin{highlight}[Reverse Algorithm]
    \begin{verbatim}
    void LinkedList::Reverse() {
        if (top_ptr_ == nullptr || top_ptr_->next == nullptr) {
            // Empty list or single node, no need to reverse
            return;
        }
        
        shared_ptr<node> previous = nullptr;
        shared_ptr<node> current = top_ptr_;
        shared_ptr<node> next = nullptr;
        
        while (current != nullptr) {
            // Find the next node 
            next = current->next;
            
            // Reverse the connection of the current node 
            current->next = previous;
            
            // Update the previous and current nodes
            previous = current;
            current = next;
        }
        // Update the top of the linked list
        top_ptr_ = previous;
    }        
    \end{verbatim}
\end{highlight}

\subsection*{Concatenating A Linked List}

Concatenating a linked list pertains to merging two linked lists. In the simplest terms, to do this, we 
simply set the tail
of a linked list to the head of the other linked list. Algorithmically, this can be seen below:

\begin{highlight}[Concatenation Algorithm]
    \begin{verbatim}
    void LinkedList::Concatenate(LinkedList& secondList) {
        // Handle the base case where the top node is null of the initial list
        if (top_ptr_ == nullptr) {
            top_ptr_ = secondList.GetTop();
        } else {
            shared_ptr<node> current = top_ptr_;
            // Find the tail of the current linked list
            while (current->next != nullptr) {
                current = current->next;
            }
            // Setting the tail of the current list to the top node of the second list
            current->next = secondList.GetTop();
        }
        // Clear the second linked list
        secondList.SetTop(nullptr);
    }        
    \end{verbatim}
\end{highlight}

\subsection*{Ascending Order In A Linked List}

To sort the data in a linked list in ascending order, we do something similar to what we do in a reverse 
algorithm in 
linked lists. We first have to create two separate nodes: `current' and `previous', and update them 
accordingly with an
if statement. This is done in the presence of a while loop that is nested in a do-while loop. 
Algorithmically, this can
be seen below.

\begin{highlight}[Ascending Order Algorithm]
    \begin{verbatim}
    void LinkedList::Sort() {
        // Handle the base case of the 
        if (top_ptr_ == nullptr || top_ptr_->next == nullptr) {
            return;
        }
        bool swapped;
        // Execute a do-while loop process to make sure the sorting
        // happens at least once
        do {
            // Set swapped to false and create new nodes
            swapped = false;
            shared_ptr<node> current = top_ptr_;
            shared_ptr<node> previous = nullptr;
            // Traverse through the list until we reach the end
            while (current->next != nullptr) {
                if (current->data > current->next->data) {
                    // Swap the data member of the current node with that of the next node
                    int temp = current->data;
                    current->data = current->next->data;
                    current->next->data = temp;
                    // Set swapped to true
                    swapped = true;
                }
                // Update the previous and current nodes
                previous = current;
                current = current->next;
            }
        } while (swapped);
    }        
    \end{verbatim}
\end{highlight}

\end{document}
% ----- ----- ----- ----- ----- End Document
